\documentclass[11pt]{article}

% ============================================================
% PACKAGES
% ============================================================

\usepackage[margin=0.8in]{geometry}
\usepackage[T1]{fontenc}
\usepackage[utf8]{inputenc}
\usepackage{lmodern}
\usepackage{microtype}

\usepackage{amsmath,amssymb}
\usepackage{array}
\usepackage{booktabs}
\usepackage{enumitem}
\usepackage{multicol}
\usepackage{ragged2e}

\usepackage{xcolor}
\usepackage[most]{tcolorbox}

\usepackage{hyperref}

% ============================================================
% HYPERREF
% ============================================================

\hypersetup{
    colorlinks=true,
    linkcolor=black,
    urlcolor=black,
    citecolor=black,
    pdftitle={History of Logic: Roman and Medieval Logic},
    pdfauthor={}
}

% ============================================================
% GENERAL FORMATTING
% ============================================================

\setlength{\parindent}{0pt}
\setlength{\parskip}{0.55em}

\setlist[itemize]{leftmargin=1.5em}
\setlist[enumerate]{leftmargin=1.7em}

% ============================================================
% CUSTOM COMMANDS
% ============================================================

\newcommand{\term}[1]{\textit{#1}}
\newcommand{\latin}[1]{\textit{#1}}
\newcommand{\thereforemark}{\(\therefore\)}
\newcommand{\imp}{\(\rightarrow\)}
\newcommand{\entails}{\(\vdash\)}

% ============================================================
% WHITEBOARD ENVIRONMENT
% ============================================================

\newtcolorbox{whiteboard}[1][]{
    breakable,
    colback=white,
    colframe=black,
    boxrule=0.8pt,
    arc=0pt,
    left=7pt,
    right=7pt,
    top=6pt,
    bottom=6pt,
    before skip=8pt,
    after skip=10pt,
    title={WHITEBOARD},
    fonttitle=\bfseries,
    #1
}

% ============================================================
% DOCUMENT
% ============================================================

\begin{document}

\begin{center}
    {\LARGE \textbf{History of Logic}}\\[0.4em]
    {\Large \textbf{Roman and Medieval Logic}}\\[0.6em]
    {\normalsize Based on William and Martha Kneale,
    \textit{The Development of Logic}, Chapter IV}
\end{center}

\vspace{0.5em}
\hrule
\vspace{0.8em}

% ============================================================
% LECTURE AIM
% ============================================================

\section*{Lecture Aim}

The Roman and medieval period is not merely an interval between
ancient logic and modern logic.

It begins with the transmission, translation, and partial fusion
of the Aristotelian and Stoic traditions, but it eventually
produces distinctively medieval theories of propositions,
semantic reference, quantification, paradox, modality, and
consequence.

The central historical movement of the lecture is:

\begin{center}
\textbf{TRANSMISSION}
\(\longrightarrow\)
\textbf{RECOVERY}
\(\longrightarrow\)
\textbf{INNOVATION}
\end{center}

By the end of the period, medieval logicians have reconstructed a
substantial body of primary, or propositional, logic and have
begun to push logical theory beyond the boundaries of traditional
Aristotelian syllogistic.

\begin{whiteboard}
\[
\boxed{
\begin{array}{c}
\text{Roman logic: preserve, translate, combine}\\
\downarrow\\
\text{Medieval logic: recover, analyze, innovate}
\end{array}
}
\]

\[
\textbf{Main themes:}
\]

\[
\begin{array}{c}
\text{language}\\
\text{meaning and reference}\\
\text{quantification}\\
\text{paradox}\\
\text{consequence}
\end{array}
\]
\end{whiteboard}

% ============================================================
% PART I
% ============================================================

\part{From Cicero to Boethius}

% ============================================================
\section{Two Ancient Traditions Become One}
% ============================================================

After Aristotle and Chrysippus, the Peripatetic and Stoic
traditions did not remain permanently separate.

Later ancient writers discussed them together, borrowed
terminology and doctrines from both, and gradually produced a
mixed logical tradition.

The result was not always a clean synthesis.

The Kneales emphasize that what developed might be described as a
fusion, but in some cases also as a confusion, of originally
distinct systems.

The historical importance of this process is enormous because the
mixed tradition becomes the logical inheritance passed to the
Latin Middle Ages.

\begin{whiteboard}
\[
\begin{array}{c}
\textbf{ARISTOTLE}\\
\text{term / syllogistic logic}
\end{array}
\qquad
+
\qquad
\begin{array}{c}
\textbf{CHRYSIPPUS}\\
\text{propositional / Stoic logic}
\end{array}
\]

\[
\downarrow
\]

\[
\text{late ancient debate and mixture}
\]

\[
\downarrow
\]

\[
\boxed{\textbf{BOETHIUS}}
\]

\[
\text{fusion --- or sometimes confusion}
\]
\end{whiteboard}

% ============================================================
\section{Cicero: Translator More Than Innovator}
% ============================================================

Cicero is not important here because he created a new formal
system.

His importance lies primarily in transmission.

Greek logical terminology had to be rendered into Latin.

Cicero helped create a Latin philosophical vocabulary that later
thinkers could inherit.

This matters because medieval logical work in western Europe
would be conducted primarily in Latin.

One important example is the development of the word
\latin{propositio}.

Its later medieval meaning as a proposition or statement was not
simply fixed from the beginning; Latin technical vocabulary
developed historically.

\begin{whiteboard}
\[
\textbf{CICERO}
\]

\[
\text{Greek logical vocabulary}
\longrightarrow
\text{Latin logical vocabulary}
\]

\[
\downarrow
\]

\[
\text{medieval logical language}
\]

\[
\boxed{\text{Cicero: TRANSMISSION $>$ INNOVATION}}
\]
\end{whiteboard}

% ============================================================
\section{Topics and Formal Validity}
% ============================================================

Cicero's treatment of topics preserves an older distinction
between two questions.

The first question is practical:

\begin{center}
Where can I \emph{find} an argument?
\end{center}

The second is logical:

\begin{center}
Does the argument actually follow?
\end{center}

Topics are associated primarily with finding appropriate lines of
argument.

Formal dialectic is concerned with whether a proposed inference
is valid.

Cicero's interests remain closely connected with rhetoric and
practical argumentation.

This reminds us that ancient and Roman logic is not yet always
separated sharply from the art of persuasive discourse.

\begin{whiteboard}
\[
\textbf{TWO QUESTIONS}
\]

\[
\begin{array}{ll}
1. & \text{How do we FIND an argument?}\\
   & \rightarrow \text{TOPICS}\\[0.6em]
2. & \text{Is the argument VALID?}\\
   & \rightarrow \text{DIALECTIC}
\end{array}
\]

\[
\boxed{\text{Cicero emphasizes argument-discovery.}}
\]
\end{whiteboard}

% ============================================================
\section{The Survival of Stoic Inferential Logic}
% ============================================================

The Stoic tradition did not disappear after Chrysippus.

Later ancient writers preserve schemes of hypothetical inference
derived from, or related to, the earlier Stoic indemonstrables.

The evidence is imperfect, but a later tradition appears to have
expanded and rearranged the basic Chrysippean material into a
larger set of hypothetical rules.

The important historical point is not the exact reconstruction of
every rule.

It is that propositional patterns of inference continued to
circulate alongside Aristotelian syllogistic.

\begin{whiteboard}
\[
\text{Chrysippus}
\]

\[
5\ \text{basic indemonstrables}
\]

\[
\downarrow
\]

\[
\text{later expansion / rearrangement}
\]

\[
\downarrow
\]

\[
7\ \text{hypothetical rules}
\]

\[
\boxed{\text{Stoic primary logic survives.}}
\]
\end{whiteboard}

% ============================================================
\section{Galen: Two Logical Traditions}
% ============================================================

Galen explicitly works with material from both the Aristotelian
and Stoic traditions.

Rather than treating one tradition as simply replacing the
other, he recognizes that they have different characteristic
applications.

Categorical syllogistic is naturally associated with geometrical
reasoning.

Hypothetical reasoning is especially useful in metaphysical and
dialectical contexts.

This is an important historical moment because the two systems are
no longer merely rivals.

They can be understood as complementary logical tools.

\begin{whiteboard}
\[
\textbf{GALEN}
\]

\[
\begin{array}{c}
\text{CATEGORICAL LOGIC}\\
\downarrow\\
\text{geometry}
\end{array}
\qquad
\begin{array}{c}
\text{HYPOTHETICAL LOGIC}\\
\downarrow\\
\text{metaphysical dialectic}
\end{array}
\]

\[
\boxed{\text{We need BOTH kinds of reasoning.}}
\]
\end{whiteboard}

% ============================================================
\section{Relational Reasoning: A Road Not Yet Taken}
% ============================================================

Galen also notices arguments whose validity depends on relations
between objects.

Consider:

\[
\text{Theon has twice as much as Dio.}
\]

\[
\text{Philo has twice as much as Theon.}
\]

\[
\therefore \text{Philo has four times as much as Dio.}
\]

Or:

\[
\text{Sophroniscus is father of Socrates.}
\]

\[
\therefore \text{Socrates is son of Sophroniscus.}
\]

These inferences cannot be represented adequately merely by
treating whole propositions as indivisible units, and they do not
fit neatly into traditional subject--predicate syllogistic.

Their validity depends upon the internal logical structure of
relations.

Galen sees the problem, but no systematic relational logic
follows from the insight.

\begin{whiteboard}
\[
A = 2B
\]

\[
C = 2A
\]

\[
\therefore C = 4B
\]

\[
\text{and}
\]

\[
X\text{ is father of }Y
\]

\[
\therefore Y\text{ is son of }X
\]

\[
\boxed{
\begin{array}{c}
\text{RELATIONAL reasoning}\\
\text{does not fit neatly into}\\
\text{ordinary syllogistic or basic Stoic forms}
\end{array}
}
\]
\end{whiteboard}

% ============================================================
\section{Porphyry and the Five Predicables}
% ============================================================

Porphyry's \textit{Isagoge} becomes one of the central works of
medieval logical education.

The famous five predicables are:

\begin{enumerate}
    \item genus,
    \item species,
    \item differentia,
    \item property,
    \item accident.
\end{enumerate}

These become known as the \latin{quinque voces}.

They provide a framework for asking how predicates relate to
subjects and how general concepts are organized.

\begin{whiteboard}
\[
\textbf{PORPHYRY: QUINQUE VOCES}
\]

\[
\begin{array}{c}
\text{GENUS}\\
\text{SPECIES}\\
\text{DIFFERENTIA}\\
\text{PROPERTY}\\
\text{ACCIDENT}
\end{array}
\]

\[
\text{Aristotle}
\longrightarrow
\text{Porphyry}
\longrightarrow
\text{medieval curriculum}
\]
\end{whiteboard}

% ============================================================
\section{Augustine and the \textit{Dicibile}}
% ============================================================

Augustine's logical reflections distinguish among words, thoughts,
and things.

The term \latin{dicibile} is used for what is understood through
a word and exists as intelligible content in the mind.

This is historically suggestive because it recalls the Stoic
distinction between a linguistic expression and what that
expression signifies.

The ancient problem of propositional content therefore remains
alive, even though the technical vocabulary has changed.

\begin{whiteboard}
\[
\textbf{AUGUSTINE}
\]

\[
\text{WORD}
\]

\[
\downarrow
\]

\[
\text{\latin{DICIBILE}}
\]

\[
\text{``what is understood''}
\]

\[
\downarrow
\]

\[
\text{thought / content}
\]

\[
\boxed{\text{possible bridge back toward the Stoic \term{lekton}}}
\]
\end{whiteboard}

% ============================================================
\section{Boethius: The Great Bridge}
% ============================================================

Boethius is the central transmitter of ancient logic to the early
medieval Latin world.

His logical translations and commentaries preserve material from
Aristotle, Porphyry, and the later ancient logical tradition.

For long periods, western students had access principally to works
such as:

\begin{itemize}
    \item Aristotle's \textit{Categories},
    \item Aristotle's \textit{De Interpretatione},
    \item Porphyry's \textit{Isagoge},
    \item Boethius' own logical works and commentaries.
\end{itemize}

Boethius is therefore not merely one author among many.

He is a historical conduit between Greek antiquity and Latin
medieval thought.

\begin{whiteboard}
\[
\textbf{BOETHIUS}
\]

\[
\text{GREEK ANTIQUITY}
\]

\[
\downarrow
\]

\[
\text{translations + commentaries}
\]

\[
\downarrow
\]

\[
\text{LATIN MIDDLE AGES}
\]

\[
\boxed{\text{BOETHIUS = THE BRIDGE}}
\]
\end{whiteboard}

% ============================================================
\section{Boethius and Hypothetical Inference}
% ============================================================

Boethius helps preserve the tradition of hypothetical logic.

A representative pattern is:

\[
A \rightarrow B,
\]

\[
B \rightarrow C,
\]

therefore:

\[
A \rightarrow C.
\]

The point is historically important because the medieval West
inherits not only categorical syllogistic but also remnants of
propositional inference.

\begin{whiteboard}
\[
\textbf{HYPOTHETICAL TRANSITIVITY}
\]

\[
\begin{array}{c}
A \rightarrow B\\
B \rightarrow C\\
\hline
A \rightarrow C
\end{array}
\]

\[
\boxed{\text{propositional material survives in Latin logic}}
\]
\end{whiteboard}

% ============================================================
\section{Boethius and \textit{Consequentia}}
% ============================================================

Boethius distinguishes between conditionals that hold merely
accidentally and those that involve a more natural or necessary
connection.

This helps preserve an important question:

\begin{center}
What does it mean for one thing to \emph{follow from} another?
\end{center}

The Latin vocabulary of \latin{consequentia} will become
increasingly important in medieval logic.

\begin{whiteboard}
\[
\textbf{CONSEQUENTIA}
\]

\[
\begin{array}{c}
\text{ACCIDENTAL}\\[0.3em]
\text{vs.}\\[0.3em]
\text{NECESSARY / NATURAL}
\end{array}
\]

\[
\boxed{\text{``following from'' becomes a central medieval problem}}
\]
\end{whiteboard}

% ============================================================
\section{Written, Spoken, and Mental Discourse}
% ============================================================

The later ancient tradition distinguishes three kinds of
discourse:

\begin{enumerate}
    \item written,
    \item spoken,
    \item mental.
\end{enumerate}

Written signs signify spoken signs.

Spoken expressions signify mental conceptions.

Mental discourse is treated as more fundamental and naturally
connected with things.

This scheme becomes extremely important in medieval theories of
mental language.

\begin{whiteboard}
\[
\textbf{THREE KINDS OF DISCOURSE}
\]

\[
\text{WRITTEN}
\]

\[
\downarrow
\]

\[
\text{SPOKEN}
\]

\[
\downarrow
\]

\[
\text{MENTAL}
\]

\[
\downarrow
\]

\[
\text{THINGS}
\]

\[
\boxed{\text{logic studies signs and what signs signify}}
\]
\end{whiteboard}

% ============================================================
\section{The Problem of Universals}
% ============================================================

Porphyry raises questions about genera and species without
resolving them.

Are universals:

\begin{itemize}
    \item realities independent of thought?
    \item objects existing only in thought?
    \item corporeal or incorporeal?
    \item separate from sensible particulars or present only in them?
\end{itemize}

Boethius transmits these problems to medieval philosophy.

The problem of universals therefore becomes one of the major
points at which logic, semantics, metaphysics, and epistemology
intersect.

\begin{whiteboard}
\[
\textbf{THE UNIVERSALS PROBLEM}
\]

\[
\begin{array}{c}
\text{GENUS / SPECIES}\\[0.4em]
\text{things?}\\
\text{thoughts?}\\
\text{words?}\\
\text{incorporeal realities?}
\end{array}
\]

\[
\text{Porphyry asks}
\longrightarrow
\text{Boethius transmits}
\longrightarrow
\text{Middle Ages debate}
\]
\end{whiteboard}

% ============================================================
% PART II
% ============================================================

\part{From Alcuin to Abelard}

% ============================================================
\section{The Early Medieval Logical Library}
% ============================================================

The early medieval West does not initially possess the whole
Aristotelian logical corpus.

The principal inherited materials include:

\begin{itemize}
    \item the \textit{Categories},
    \item \textit{De Interpretatione},
    \item Porphyry's \textit{Isagoge},
    \item the logical writings of Boethius.
\end{itemize}

These works later become associated with the \latin{ars vetus},
the ``old logic.''

The limited textual inheritance is important because medieval
thinkers begin rebuilding logical study from a comparatively
narrow collection.

\begin{whiteboard}
\[
\textbf{EARLY MEDIEVAL LOGIC}
\]

\[
\begin{array}{c}
\textit{Categories}\\
\textit{De Interpretatione}\\
\textit{Isagoge}\\
\text{Boethius}
\end{array}
\]

\[
\downarrow
\]

\[
\boxed{\text{\latin{ARS VETUS}}}
\]
\end{whiteboard}

% ============================================================
\section{Alcuin and the Trivium}
% ============================================================

Alcuin is important less for original logical theory than for
education.

The traditional trivium is:

\begin{enumerate}
    \item grammar,
    \item dialectic,
    \item rhetoric.
\end{enumerate}

Dialectic thereby becomes one of the basic disciplines of higher
intellectual training.

\begin{whiteboard}
\[
\textbf{TRIVIUM}
\]

\[
\begin{array}{c}
\text{GRAMMAR}\\
\text{DIALECTIC}\\
\text{RHETORIC}
\end{array}
\]

\[
\boxed{\text{logic becomes structurally embedded in education}}
\]
\end{whiteboard}

% ============================================================
\section{Grammar and Logic}
% ============================================================

The recovery of logic occurs alongside intense study of grammar.

This creates both opportunities and confusions.

Medieval thinkers repeatedly ask how to distinguish:

\begin{itemize}
    \item the physical or written word,
    \item the meaning of the word,
    \item the thing named,
    \item the universal associated with the term.
\end{itemize}

These questions gradually become logical and semantic problems in
their own right.

\begin{whiteboard}
\[
\textbf{MEDIEVAL SEMANTIC PROBLEM}
\]

\[
\text{WORD}
\]

\[
?
\]

\[
\text{MEANING}
\]

\[
?
\]

\[
\text{THING}
\]

\[
\boxed{
\begin{array}{c}
\text{Grammar + Logic}\\
\text{can illuminate one another}\\
\text{or confuse one another}
\end{array}
}
\]
\end{whiteboard}

% ============================================================
\section{Nominalism and Realism}
% ============================================================

The proposition

\[
\text{\latin{Homo est species}}
\]

means:

\begin{center}
Man is a species.
\end{center}

But what does the subject term stand for?

If a species is something real, we move toward realism.

If the universal is only a word or sound, we move toward
nominalism.

Roscellin becomes associated with a strong nominalist tendency
according to which universals are merely words or a
\latin{flatus vocis}.

Abelard rejects simple realism but also does not accept a crude
theory in which universals are nothing but noises.

\begin{whiteboard}
\[
\textbf{UNIVERSALS}
\]

\[
\begin{array}{c}
\text{REALISM}\\
\text{universal has reality}
\end{array}
\]

\[
\text{vs.}
\]

\[
\begin{array}{c}
\text{NOMINALISM}\\
\text{universal = word / \latin{flatus vocis}}
\end{array}
\]

\[
\boxed{\text{ABELARD: neither simple extreme}}
\]
\end{whiteboard}

% ============================================================
\section{A New Intellectual Confidence}
% ============================================================

A major transformation occurs in the twelfth century.

The goal is no longer merely to recover and preserve the ancients.

Thinkers increasingly assume that ancient philosophy can be
analyzed, criticized, and improved.

This change in intellectual confidence is essential for
understanding Abelard.

\begin{whiteboard}
\[
\text{EARLY MIDDLE AGES}
\]

\[
\text{recover the ancients}
\]

\[
\downarrow
\]

\[
\textbf{TWELFTH CENTURY}
\]

\[
\text{argue with the ancients}
\]

\[
\downarrow
\]

\[
\boxed{\text{improve the ancients}}
\]
\end{whiteboard}

% ============================================================
\section{Anselm and Dialectical Theology}
% ============================================================

Anselm applies dialectical reasoning directly to theological
questions.

One important method is reductive reasoning.

Suppose the denial of a proposition.

If that supposition produces an impossibility or contradiction,
reject the denial.

Schematically:

\[
\neg P \rightarrow \bot
\]

therefore:

\[
P.
\]

The important point for the history of logic is methodological:
logical patterns inherited from antiquity are now being used
systematically in theological argument.

\begin{whiteboard}
\[
\textbf{ANSELM}
\]

\[
\text{suppose }\neg P
\]

\[
\downarrow
\]

\[
\text{derive contradiction}
\]

\[
\downarrow
\]

\[
\boxed{P}
\]

\[
\text{ancient dialectical forms}
\rightarrow
\text{medieval theology}
\]
\end{whiteboard}

% ============================================================
\section{The Scholastic \textit{Quaestio}}
% ============================================================

Abelard's \textit{Sic et Non} exemplifies a style of inquiry that
becomes characteristic of scholastic intellectual life.

One begins with apparently conflicting claims or authorities.

The task is not simply to choose one side immediately.

Instead one must identify distinctions that allow the apparent
contradiction to be analyzed.

\begin{whiteboard}
\[
\textbf{SCHOLASTIC QUAESTIO}
\]

\[
\text{PRO}
\]

\[
\text{vs.}
\]

\[
\text{CONTRA}
\]

\[
\downarrow
\]

\[
\text{DISTINCTIONS}
\]

\[
\downarrow
\]

\[
\boxed{\text{RESOLUTION}}
\]

\[
\text{Reason advances through organized disputation.}
\]
\end{whiteboard}

% ============================================================
\section{Abelard as an Innovator}
% ============================================================

Abelard is not merely a compiler.

His logical work shows confidence that the ancient tradition is
incomplete and that further discoveries can be made.

His \textit{Dialectica} therefore marks an important transition
from recovery to active development.

\begin{whiteboard}
\[
\textbf{ABELARD}
\]

\[
\boxed{\text{Ancient logic} \neq \text{finished logic}}
\]

\[
\text{Logic remains a field of discovery.}
\]
\end{whiteboard}

% ============================================================
\section{What Is True or False?}
% ============================================================

Abelard defines a proposition as discourse signifying truth or
falsity.

Yet he recognizes that the fundamental concern of logic cannot be
merely the physical utterance or written sentence.

A proposition is logically interesting because of what it puts
forward as being the case.

Thus we need to distinguish the propositional sign from what is
expressed by the sign.

\begin{whiteboard}
\[
\textbf{ABELARD}
\]

\[
\text{PROPOSITION-SIGN}
\]

\[
\downarrow
\]

\[
\text{expresses}
\]

\[
\downarrow
\]

\[
\text{WHAT IS THE CASE}
\]

\[
\downarrow
\]

\[
\text{TRUE / FALSE}
\]

\[
\boxed{\text{truth concerns propositional content}}
\]
\end{whiteboard}

% ============================================================
\section{The \textit{Dictum}}
% ============================================================

Abelard uses the term \latin{dictum} for what is expressed by a
proposition.

For example:

\[
\text{Socrates is a man}
\]

expresses:

\[
\text{Socrates' being-a-man.}
\]

But the \latin{dictum} is not simply another ordinary object
alongside Socrates and other things.

The distinction is one of the important medieval attempts to
understand propositional content.

\begin{whiteboard}
\[
\text{``Socrates is a man''}
\]

\[
\downarrow
\]

\[
\textbf{\latin{DICTUM}}
\]

\[
\text{Socrates' being-a-man}
\]

\[
\boxed{\text{what is expressed} \neq \text{another ordinary thing}}
\]
\end{whiteboard}

% ============================================================
\section{The Copula}
% ============================================================

Abelard gives special logical attention to the verb \latin{est}.

In a proposition such as:

\[
\text{Socrates is a man}
\]

the word ``is'' functions as a copula joining subject and
predicate.

It should not automatically be understood as an independent
assertion of existence.

This distinction is important because otherwise every use of the
copula would appear to commit the speaker to the existence of the
subject in a way that ordinary predication does not.

\begin{whiteboard}
\[
\text{Socrates \textbf{IS} a man}
\]

\[
\hspace{0.4in}\uparrow
\]

\[
\textbf{COPULA}
\]

\[
\boxed{\text{copula} \neq \text{independent assertion of existence}}
\]
\end{whiteboard}

% ============================================================
\section{Toward Supposition Theory}
% ============================================================

A central medieval question becomes:

\begin{center}
What does a term stand for in a proposition?
\end{center}

Abelard's semantic analysis helps prepare the way for the later
theory of \latin{suppositio terminorum}.

The same word can perform different logical roles in different
contexts.

The later medievals will turn this insight into an elaborate
technical theory.

\begin{whiteboard}
\[
\textbf{TERM IN A PROPOSITION}
\]

\[
\boxed{\text{What does it stand for HERE?}}
\]

\[
\text{ABELARD}
\]

\[
\downarrow
\]

\[
\text{\latin{SUPPOSITIO TERMINORUM}}
\]
\end{whiteboard}

% ============================================================
\section{Negation Has Scope}
% ============================================================

Abelard sees that negation can apply to a whole proposition rather
than merely to one term inside it.

Compare:

\[
\text{Not every man is white}
\]

with:

\[
\text{Some man is not white.}
\]

The two should not simply be treated as identical without
examining assumptions about existence and quantification.

This is a major step toward explicit analysis of logical scope.

\begin{whiteboard}
\[
\textbf{SCOPE OF NEGATION}
\]

\[
\neg[\text{Every man is white}]
\]

\[
\neq \text{ automatically}
\]

\[
\text{Some man is }\neg\text{ white}
\]

\[
\boxed{
\begin{array}{c}
\text{negation of the WHOLE}\\
\text{must be distinguished from}\\
\text{negation INSIDE the proposition}
\end{array}
}
\]
\end{whiteboard}

% ============================================================
\section{Existential Import and the Square}
% ============================================================

Abelard worries about the traditional square of opposition
because he takes universal affirmative propositions to carry
existential import.

Consider:

\[
\text{Every man is a man.}
\]

If there are no men, should this still count as true?

Abelard's answer is no.

This shows that the traditional relations among quantified
propositions contain assumptions about existence that must be
made explicit.

\begin{whiteboard}
\[
\textbf{EXISTENTIAL IMPORT}
\]

\[
\text{Every man is a man.}
\]

\[
\text{Suppose: NO MEN EXIST}
\]

\[
\text{True or false?}
\]

\[
\boxed{\text{ABELARD: FALSE}}
\]

\[
\text{The traditional square contains existential assumptions.}
\]
\end{whiteboard}

% ============================================================
\section{Quantifiers as Logical Operators}
% ============================================================

Abelard explicitly analyzes the signs of quantity:

\[
\latin{omnis},
\qquad
\latin{quidam},
\qquad
\latin{nullus}.
\]

That is:

\[
\text{every},
\qquad
\text{some},
\qquad
\text{no}.
\]

The logical form of a proposition depends in part on how these
expressions determine the range of the subject term.

\begin{whiteboard}
\[
\textbf{SIGNA QUANTITATIS}
\]

\[
\begin{array}{rcl}
\latin{OMNIS} &\rightarrow& \text{EVERY}\\
\latin{QUIDAM} &\rightarrow& \text{SOME}\\
\latin{NULLUS} &\rightarrow& \text{NO}
\end{array}
\]

\[
\boxed{\text{logical form depends on quantifier behavior}}
\]
\end{whiteboard}

% ============================================================
\section{Modal Scope}
% ============================================================

Abelard also distinguishes two ways modal language can operate.

Compare:

\[
\text{It is possible that no man is white}
\]

with:

\[
\text{No man is possibly white.}
\]

In the first case, possibility applies to an entire proposition.

In the second, the modal expression appears within the
predication.

This anticipates the later distinction between modality
\latin{de dicto} and modality \latin{de re}.

\begin{whiteboard}
\[
\textbf{MODAL SCOPE}
\]

\[
\text{POSSIBLE }[\text{No man is white}]
\]

\[
\neq
\]

\[
\text{No man is }[\text{POSSIBLY white}]
\]

\[
\boxed{
\begin{array}{c}
\text{later terminology:}\\
\latin{DE DICTO}\quad \text{vs.}\quad \latin{DE RE}
\end{array}
}
\]
\end{whiteboard}

% ============================================================
\section{Abelard's Theory of Consequence}
% ============================================================

For Abelard, a satisfactory conditional expresses a necessary
connection.

He speaks of \latin{necessitas consecutionis}.

He distinguishes arguments whose validity is secured by their
formal arrangement from arguments that require information about
the particular things denoted by their terms.

The contrast can be expressed as structure versus special facts
about the terms.

\begin{whiteboard}
\[
\textbf{ABELARD: CONSEQUENCE}
\]

\[
\begin{array}{c}
\text{FORMAL / PERFECT}\\
\text{truth from STRUCTURE}
\end{array}
\]

\[
\text{vs.}
\]

\[
\begin{array}{c}
\text{IMPERFECT}\\
\text{needs facts about TERMS / THINGS}
\end{array}
\]
\end{whiteboard}

% ============================================================
\section{The Paradox of Implication}
% ============================================================

Suppose the antecedent of a conditional is impossible.

For example:

\[
\text{Socrates is a stone.}
\]

Then the combination:

\[
\text{Socrates is a stone and Socrates is not an ass}
\]

is itself impossible.

This makes it tempting to say:

\[
\text{If Socrates is a stone, then Socrates is an ass.}
\]

But Abelard resists treating such a conditional as expressing a
genuine intrinsic connection.

For strict consequence, the antecedent should require the
consequent \latin{ex se ipso}, from itself.

\begin{whiteboard}
\[
\textbf{PARADOX OF IMPLICATION}
\]

\[
\text{Impossible }P
\]

\[
\text{Does }P \rightarrow Q
\]

\[
\text{hold for ANY }Q?
\]

\[
\boxed{
\begin{array}{c}
\text{ABELARD: not in the strict sense}\\
\text{Antecedent must require consequent intrinsically}
\end{array}
}
\]
\end{whiteboard}

% ============================================================
\section{Abelard's Inferential Rules}
% ============================================================

Abelard explicitly works with several familiar patterns.

\subsection*{Modus Ponens}

\[
P \rightarrow Q
\]

\[
P
\]

\[
\therefore Q
\]

\subsection*{Modus Tollens}

\[
P \rightarrow Q
\]

\[
\neg Q
\]

\[
\therefore \neg P
\]

\subsection*{Transitivity}

\[
P \rightarrow Q
\]

\[
Q \rightarrow R
\]

\[
\therefore P \rightarrow R
\]

\begin{whiteboard}
\[
\textbf{ABELARD}
\]

\[
\frac{P\rightarrow Q \qquad P}{Q}
\]

\[
\frac{P\rightarrow Q \qquad \neg Q}{\neg P}
\]

\[
\frac{P\rightarrow Q \qquad Q\rightarrow R}{P\rightarrow R}
\]

\[
\boxed{
\begin{array}{c}
\text{modus ponens}\\
\text{modus tollens}\\
\text{transitivity}
\end{array}
}
\]
\end{whiteboard}

% ============================================================
% PART III
% ============================================================

\part{The Logic of the Universities}

% ============================================================
\section{The \textit{Ars Nova}}
% ============================================================

By the middle of the twelfth century, much more of Aristotle's
logical corpus becomes available in Latin.

The newly recovered portions of the \textit{Organon} become known
as the \latin{ars nova}.

The older core remains the \latin{ars vetus}.

The result is a greatly expanded logical curriculum.

\begin{whiteboard}
\[
\textbf{\latin{ARS VETUS}}
\]

\[
+
\]

\[
\textbf{\latin{ARS NOVA}}
\]

\[
\downarrow
\]

\[
\boxed{\text{UNIVERSITY LOGIC}}
\]
\end{whiteboard}

% ============================================================
\section{Logic Becomes a University Discipline}
% ============================================================

Logic becomes central to the arts curriculum at universities such
as Paris and Oxford.

Its characteristic forms include:

\begin{itemize}
    \item lectures,
    \item commentaries,
    \item \latin{quaestiones},
    \item disputations,
    \item handbooks and \latin{summulae}.
\end{itemize}

Logical reasoning is no longer merely something read in ancient
texts.

It is a practiced academic discipline.

\begin{whiteboard}
\[
\textbf{MEDIEVAL UNIVERSITY LOGIC}
\]

\[
\text{LECTURE}
\]

\[
\downarrow
\]

\[
\text{\latin{QUAESTIO}}
\]

\[
\downarrow
\]

\[
\text{DISPUTATION}
\]

\[
\downarrow
\]

\[
\boxed{\text{FORMAL TRAINING}}
\]
\end{whiteboard}

% ============================================================
\section{Sophisms as Logical Laboratories}
% ============================================================

The recovery of Aristotle's \textit{De Sophisticis Elenchis}
stimulates an enormous interest in sophisms.

A \latin{sophisma} can be used as a deliberately difficult case.

The point is not merely to avoid mistakes.

The puzzle forces the logician to identify ambiguity, scope,
reference, and inference patterns.

Sophisms therefore function as intellectual laboratories.

\begin{whiteboard}
\[
\textbf{SOPHISMA}
\]

\[
\text{construct a difficult case}
\]

\[
\downarrow
\]

\[
\text{force distinctions}
\]

\[
\downarrow
\]

\[
\text{test logical theory}
\]

\[
\boxed{\text{logical puzzle = laboratory for analysis}}
\]
\end{whiteboard}

% ============================================================
\section{\textit{Insolubilia}: The Liar Returns}
% ============================================================

Variants of the Liar paradox become known as
\latin{insolubilia}.

A simple example is:

\[
\text{``This proposition is false.''}
\]

The medievals recognize that the difficulty involves a peculiar
kind of self-reference.

Different solutions are proposed, including:

\begin{itemize}
    \item \latin{cassatio},
    \item \latin{restrictio},
    \item solution \latin{secundum quid et simpliciter}.
\end{itemize}

\latin{Cassatio} attempts to deny that the problematic expression
successfully states anything.

\latin{Restrictio} restricts the range of reference in order to
block the vicious self-reference.

\begin{whiteboard}
\[
\textbf{\latin{INSOLUBILIA}}
\]

\[
\text{``This proposition is false.''}
\]

\[
\begin{array}{c}
\latin{CASSATIO}\\
\latin{RESTRICTIO}\\
\latin{SECUNDUM\ QUID\ /\ SIMPLICITER}
\end{array}
\]

\[
\boxed{\text{self-reference becomes a technical logical problem}}
\]
\end{whiteboard}

% ============================================================
\section{\textit{Syncategoremata}}
% ============================================================

Medieval logicians write special treatises on expressions that do
not ordinarily function as terms but are essential to logical
structure.

Examples include:

\[
\text{and,\ or,\ not,\ if,\ every,\ some,\ only,\ except.}
\]

These are \latin{syncategoremata}.

The important insight is that logical form can be controlled by
very small pieces of language.

\begin{whiteboard}
\[
\textbf{\latin{SYNCATEGOREMATA}}
\]

\[
\begin{array}{c}
\text{AND}\\
\text{OR}\\
\text{NOT}\\
\text{IF}\\
\text{EVERY}\\
\text{SOME}\\
\text{ONLY}\\
\text{EXCEPT}
\end{array}
\]

\[
\boxed{\text{small words --- large effects on LOGICAL FORM}}
\]
\end{whiteboard}

% ============================================================
\section{\textit{Obligationes}}
% ============================================================

An \latin{obligatio} is a regulated disputational situation.

A participant agrees to an assumption or set of conditions and
must then answer subsequent questions consistently within that
framework.

This turns logical disputation into a rule-governed exercise.

It continues the ancient interest in dialectical debate but with a
more elaborate formal structure.

\begin{whiteboard}
\[
\textbf{\latin{OBLIGATIONES}}
\]

\[
\text{assume }P
\]

\[
\downarrow
\]

\[
\text{conduct dispute}
\]

\[
\downarrow
\]

\[
\text{maintain consistency}
\]

\[
\boxed{\text{logic as a rule-governed dialectical game}}
\]
\end{whiteboard}

% ============================================================
\section{William of Shyreswood and Peter of Spain}
% ============================================================

William of Shyreswood's \textit{Introductiones in Logicam}
organizes material on:

\begin{itemize}
    \item propositions,
    \item predicables,
    \item syllogisms,
    \item dialectical topics,
    \item properties of terms,
    \item fallacies.
\end{itemize}

Peter of Spain's \textit{Summulae Logicales} becomes an even more
influential logical handbook.

The university manuals combine inherited Aristotelian doctrine
with specifically medieval developments.

\begin{whiteboard}
\[
\textbf{HIGH MEDIEVAL MANUALS}
\]

\[
\text{William of Shyreswood}
\]

\[
\downarrow
\]

\[
\text{Peter of Spain}
\]

\[
\downarrow
\]

\[
\text{\latin{SUMMULAE} tradition}
\]

\[
\boxed{\text{ancient logic + medieval innovations}}
\]
\end{whiteboard}

% ============================================================
\section{Mental Language}
% ============================================================

The theory of written, spoken, and mental discourse receives new
development under the influence of Aristotle and Arabic
philosophy.

Mental terms are treated as \latin{intentiones}, natural signs or
concepts in the soul.

This creates a theory in which conventional spoken and written
signs are subordinated to a more fundamental mental language.

\begin{whiteboard}
\[
\textbf{MENTAL LANGUAGE}
\]

\[
\text{WRITTEN SIGNS}
\]

\[
\downarrow
\]

\[
\text{SPOKEN SIGNS}
\]

\[
\downarrow
\]

\[
\text{MENTAL SIGNS}
\]

\[
\downarrow
\]

\[
\text{THINGS}
\]

\[
\boxed{\latin{intentio} = \text{mental sign / concept}}
\]
\end{whiteboard}

% ============================================================
\section{Ockham's Systematic Logic}
% ============================================================

William Ockham attempts to organize the entire subject of logic.

His \textit{Summa Totius Logicae} is divided broadly into:

\begin{enumerate}
    \item terms,
    \item propositions,
    \item arguments.
\end{enumerate}

Within this framework he incorporates both traditional logical
material and specifically medieval developments.

\begin{whiteboard}
\[
\textbf{OCKHAM}
\]

\[
\begin{array}{lll}
\text{BOOK I} &\rightarrow& \text{TERMS}\\
\text{BOOK II} &\rightarrow& \text{PROPOSITIONS}\\
\text{BOOK III} &\rightarrow& \text{ARGUMENTS}
\end{array}
\]

\[
\boxed{\text{ancient inheritance + medieval logical theory}}
\]
\end{whiteboard}

% ============================================================
\section{The Fourteenth-Century Logical Culture}
% ============================================================

After Ockham, important figures include:

\begin{itemize}
    \item Walter Burleigh,
    \item Jean Buridan,
    \item Richard Swineshead,
    \item Albert of Saxony,
    \item Marsilius of Inghen,
    \item William Heytesbury,
    \item Ralph Strode,
    \item Richard Ferrybridge.
\end{itemize}

This is a period of exceptionally intense logical activity.

Medieval thinkers continue refining the theories of terms,
consequences, modality, sophisms, paradox, and disputation.

\begin{whiteboard}
\[
\textbf{FOURTEENTH-CENTURY LOGIC}
\]

\[
\begin{array}{c}
\text{Burleigh}\\
\text{Buridan}\\
\text{Swineshead}\\
\text{Albert of Saxony}\\
\text{Heytesbury}\\
\text{Strode}\\
\text{others}
\end{array}
\]

\[
\boxed{\text{one of the most intensive periods of logical study}}
\]
\end{whiteboard}

% ============================================================
\section{\textit{Via Antiqua} and \textit{Via Moderna}}
% ============================================================

The later medieval intellectual world develops competing
educational tendencies.

The \latin{via antiqua} is associated with older Aristotelian and
realist approaches.

The \latin{via moderna} is associated more strongly with terminist
and nominalist developments and with newer kinds of logical
analysis.

The difference eventually becomes institutional.

\begin{whiteboard}
\[
\textbf{\latin{VIA ANTIQUA}}
\]

\[
\text{realism + older Aristotelian emphasis}
\]

\[
\text{vs.}
\]

\[
\textbf{\latin{VIA MODERNA}}
\]

\[
\text{terminism / nominalist tendency}
\]

\[
\text{newer logical analysis}
\]
\end{whiteboard}

% ============================================================
% PART IV
% ============================================================

\part{\textit{Proprietates Terminorum}}

% ============================================================
\section{The Properties of Terms}
% ============================================================

The theory of the \latin{proprietates terminorum} develops in the
later twelfth century and becomes one of the most distinctive
areas of medieval logic.

The basic question is:

\begin{center}
What logical role does a word or phrase perform when it occurs as
a term in a proposition?
\end{center}

William of Shyreswood gives four principal properties:

\begin{enumerate}
    \item \latin{significatio},
    \item \latin{suppositio},
    \item \latin{copulatio},
    \item \latin{appellatio}.
\end{enumerate}

\begin{whiteboard}
\[
\textbf{\latin{PROPRIETATES TERMINORUM}}
\]

\[
\begin{array}{c}
\latin{SIGNIFICATIO}\\
\latin{SUPPOSITIO}\\
\latin{COPULATIO}\\
\latin{APPELLATIO}
\end{array}
\]

\[
\boxed{\text{What is a term doing IN THIS PROPOSITION?}}
\]
\end{whiteboard}

% ============================================================
\section{\textit{Significatio} and \textit{Suppositio}}
% ============================================================

A useful introductory distinction is:

\[
\latin{significatio}
\]

concerns what a term means or signifies.

\[
\latin{suppositio}
\]

concerns what the term stands for in a particular propositional
context.

The same expression can retain a general meaning while standing
for different kinds of things in different occurrences.

\begin{whiteboard}
\[
\textbf{\latin{SIGNIFICATIO}}
\]

\[
\text{``What does the word MEAN?''}
\]

\[
\textbf{\latin{SUPPOSITIO}}
\]

\[
\text{``What does it stand for HERE?''}
\]

\[
\boxed{\text{meaning} \neq \text{contextual reference}}
\]
\end{whiteboard}

% ============================================================
\section{Three Core Types of Supposition}
% ============================================================

The full medieval taxonomies are extremely elaborate.

For introductory purposes, three central cases are especially
useful.

\subsection*{Material Supposition}

The expression stands for the linguistic sign itself.

Example:

\[
\text{``Man'' is a word.}
\]

\subsection*{Simple Supposition}

The expression stands for a universal, concept, or associated
mental sign.

Example:

\[
\text{Man is a species.}
\]

\subsection*{Personal Supposition}

The expression stands for the individual things falling under the
term.

Example:

\[
\text{Every man is mortal.}
\]

\begin{whiteboard}
\[
\textbf{THREE CORE CASES}
\]

\[
\begin{array}{ll}
1. & \text{MATERIAL}\\
   & \text{``man'' as linguistic sign}\\[0.5em]
2. & \text{SIMPLE}\\
   & \text{man as concept / universal}\\[0.5em]
3. & \text{PERSONAL}\\
   & \text{individual men}
\end{array}
\]

\[
\boxed{\text{same term --- different logical role}}
\]
\end{whiteboard}

% ============================================================
\section{Personal Supposition and Quantification}
% ============================================================

Within personal supposition, medieval thinkers distinguish
different ways in which a common term ranges over individuals.

A simplified teaching scheme is:

\begin{itemize}
    \item determinate,
    \item confused,
    \item distributive.
\end{itemize}

These distinctions function partly as a substitute for what
modern logic handles through explicit quantifier notation.

\begin{whiteboard}
\[
\textbf{PERSONAL SUPPOSITION}
\]

\[
\begin{array}{c}
\text{DETERMINATE}\\
\rightarrow \text{some particular member}
\end{array}
\]

\[
\begin{array}{c}
\text{CONFUSED}\\
\rightarrow \text{range not individually fixed}
\end{array}
\]

\[
\begin{array}{c}
\text{DISTRIBUTIVE}\\
\rightarrow \text{all relevant members}
\end{array}
\]

\[
\boxed{\text{medieval substitute for parts of quantification theory}}
\]
\end{whiteboard}

% ============================================================
\section{Ockham's Nominalist Revision}
% ============================================================

William of Shyreswood and Peter of Spain work within a framework
in which a general term can be said to signify a common form or
universal.

Ockham rejects the existence of such a universal entity.

For him, signs ultimately signify individual things.

This makes \latin{suppositio} increasingly central to semantic
theory.

\begin{whiteboard}
\[
\textbf{WILLIAM / PETER}
\]

\[
\text{general term}
\rightarrow
\text{common form / universal}
\]

\[
\text{vs.}
\]

\[
\textbf{OCKHAM}
\]

\[
\text{general term}
\rightarrow
\text{individual things}
\]

\[
\boxed{\text{\latin{suppositio} becomes more fundamental}}
\]
\end{whiteboard}

% ============================================================
\section{Burleigh's Objection}
% ============================================================

Walter Burleigh objects to Ockham's account of signification.

Suppose the word ``man'' directly meant Socrates, Plato, and the
other individual men.

Then learning the meaning of the word ``man'' should require
learning that the word applies to those individuals.

But someone can understand the word ``man'' without knowing
Socrates.

Therefore the meaning of a general term cannot simply be
identified with a list of individuals.

\begin{whiteboard}
\[
\textbf{BURLEIGH AGAINST OCKHAM}
\]

\[
\text{If ``man'' meant}
\]

\[
\text{Socrates + Plato + }\cdots
\]

\[
\text{then learning ``man''}
\]

\[
\text{would require knowing Socrates.}
\]

\[
\text{But it does not.}
\]

\[
\boxed{\text{MEANING} \neq \text{list of individuals}}
\]
\end{whiteboard}

% ============================================================
\section{Distribution}
% ============================================================

The doctrine of distribution is the most familiar surviving part
of medieval term theory.

A term is distributed when it is taken for all the individuals to
which it applies.

For example:

\[
\text{Every man is mortal.}
\]

The subject term ``man'' is taken universally.

Later traditional syllogistic develops two famous rules:

\begin{enumerate}
    \item The middle term must be distributed at least once in the premises.
    \item A term may not be distributed in the conclusion unless it was distributed in a premise.
\end{enumerate}

\begin{whiteboard}
\[
\textbf{DISTRIBUTION}
\]

\[
\text{Every MAN is mortal.}
\]

\[
\text{MAN}
\rightarrow
\text{taken for every man}
\]

\[
\boxed{
\begin{array}{l}
1.\ \text{Middle term distributed at least once.}\\
2.\ \text{No new distribution in the conclusion.}
\end{array}
}
\]
\end{whiteboard}

% ============================================================
\section{What Term Theory Was Trying to Do}
% ============================================================

The theory of the properties of terms combines several different
projects.

First, it attempts a general theory of language involving meaning,
application, and reference.

Second, it analyzes the detailed behavior of the natural language
used by philosophers.

Third, it develops rules for quantified inference.

These are all respectable projects, but they do not necessarily
belong inside one single theory.

\begin{whiteboard}
\[
\textbf{\latin{PROPRIETATES TERMINORUM}}
\]

\[
\begin{array}{ll}
1. & \text{theory of LANGUAGE}\\
   & \text{meaning / reference}\\[0.5em]
2. & \text{analysis of logical IDIOMS}\\[0.5em]
3. & \text{rules of QUANTIFICATION}
\end{array}
\]

\[
\boxed{\text{valuable projects --- but not naturally one theory}}
\]
\end{whiteboard}

% ============================================================
\section{The Ambiguity of \textit{Suppositio}}
% ============================================================

The basic distinction between signification and supposition is
useful.

The difficulty is that the word \latin{suppositio} becomes
overloaded.

Compare:

\[
\text{``Socrates'' stands for Socrates}
\]

with:

\[
\text{``man'' stands for Socrates, Plato, and others.}
\]

These are not obviously the same semantic relation.

Yet both can be described using the language of ``standing for.''

\begin{whiteboard}
\[
\textbf{PROBLEM WITH \latin{SUPPOSITIO}}
\]

\[
\text{1. A name} \rightarrow \text{Socrates}
\]

\[
\text{2. MAN} \rightarrow \text{Socrates, Plato, }\ldots
\]

\[
\text{Both called: ``stands for''}
\]

\[
\boxed{\text{but these are different semantic relations}}
\]
\end{whiteboard}

% ============================================================
% PART V
% ============================================================

\part{\textit{Consequentiae}}

% ============================================================
\section{What Is a \textit{Consequentia}?}
% ============================================================

The later medieval theory of \latin{consequentiae} develops from
earlier work associated with Abelard, although the exact
historical chain is not completely known.

The central question is:

\begin{center}
Under what conditions does one proposition validly follow from
another?
\end{center}

This broadens logical investigation beyond the traditional theory
of categorical syllogism.

\begin{whiteboard}
\[
\textbf{\latin{CONSEQUENTIA}}
\]

\[
P
\]

\[
\downarrow
\]

\[
Q
\]

\[
\boxed{\text{WHEN DOES }Q\text{ FOLLOW FROM }P\text{?}}
\]
\end{whiteboard}

% ============================================================
\section{Two Paradoxes of Consequence}
% ============================================================

Two medieval principles can appear surprising.

First:

\[
\text{anything follows from the impossible.}
\]

In modern notation:

\[
\bot \vdash P.
\]

Second:

\[
\text{the necessary follows from anything.}
\]

If \(N\) is necessary, then:

\[
P \vdash N.
\]

The intuitive explanation concerns the impossibility of a
counterexample in which the antecedent is true and the consequent
false.

\begin{whiteboard}
\[
\textbf{PARADOXES OF CONSEQUENCE}
\]

\[
\text{IMPOSSIBLE}
\]

\[
\downarrow
\]

\[
\text{ANYTHING}
\]

\[
\text{and}
\]

\[
\text{ANYTHING}
\]

\[
\downarrow
\]

\[
\text{NECESSARY TRUTH}
\]

\[
\boxed{\text{Can there be a true antecedent and false consequent?}}
\]
\end{whiteboard}

% ============================================================
\section{Formal and Material Consequence}
% ============================================================

A particularly important medieval distinction separates formal
from material consequence.

A consequence is formal when its validity remains intact under
suitable replacement of its non-logical terms while the logical
structure is preserved.

A material consequence depends upon the special meanings of the
terms or upon additional facts.

\begin{whiteboard}
\[
\textbf{FORMAL CONSEQUENCE}
\]

\[
\text{valid under replacement}
\]

\[
\text{while FORM stays fixed}
\]

\[
\text{vs.}
\]

\[
\textbf{MATERIAL CONSEQUENCE}
\]

\[
\text{valid because of particular meanings / facts}
\]
\end{whiteboard}

% ============================================================
\section{Substitution as a Test of Form}
% ============================================================

Consider:

\[
\text{Every }M\text{ is }L.
\]

\[
\text{Every }S\text{ is }M.
\]

\[
\therefore \text{Every }S\text{ is }L.
\]

Now replace \(M\), \(L\), and \(S\) with different suitable
terms.

If the argument remains valid under the substitutions, then the
validity belongs to the structure rather than the specific
subject matter.

\begin{whiteboard}
\[
\textbf{FORMAL TEST}
\]

\[
\begin{array}{c}
\text{Every }M\text{ is }L\\
\text{Every }S\text{ is }M\\
\hline
\text{Every }S\text{ is }L
\end{array}
\]

\[
\boxed{
\begin{array}{c}
\text{replace }M,L,S\\
\text{validity survives}\\
\Rightarrow \text{FORM carries the inference}
\end{array}
}
\]
\end{whiteboard}

% ============================================================
\section{\textit{Bona Simpliciter} and \textit{Bona Ut Nunc}}
% ============================================================

Medieval logicians also distinguish kinds of material consequence.

A consequence can be good \latin{simpliciter}, because a stable
background relation between the terms supplies what is needed.

Or it can be good \latin{ut nunc}, ``as things now stand,'' because
some contingent circumstance currently makes the inference work.

Suppose:

\[
\text{Socrates is white.}
\]

Then:

\[
\text{Socrates runs}
\]

may support:

\[
\text{A white thing runs.}
\]

But this works because of the contingent fact that Socrates is
white.

\begin{whiteboard}
\[
\textbf{TWO MATERIAL CASES}
\]

\[
\begin{array}{c}
\latin{BONA SIMPLICITER}\\
\text{stable background truth}
\end{array}
\]

\[
\text{vs.}
\]

\[
\begin{array}{c}
\latin{BONA UT NUNC}\\
\text{true given present circumstances}
\end{array}
\]

\[
\text{Socrates is white}
\]

\[
\text{Socrates runs}
\]

\[
\therefore \text{A white thing runs}
\]
\end{whiteboard}

% ============================================================
\section{Ockham's General Principles}
% ============================================================

Ockham gives a large set of rules governing consequence.

For lecture purposes, the following family captures the central
ideas:

\begin{enumerate}
    \item A true proposition does not validly entail a false proposition.
    \item If \(P\) entails \(Q\), then the negation of \(Q\) entails the negation of \(P\).
    \item If \(P\) entails \(Q\), and \(Q\) entails \(R\), then \(P\) entails \(R\).
    \item Anything follows from the impossible.
    \item A necessary truth follows from anything.
\end{enumerate}

\begin{whiteboard}
\[
\textbf{OCKHAM: GENERAL PRINCIPLES}
\]

\[
\begin{array}{l}
1.\ \text{TRUE does not entail FALSE}\\[0.4em]
2.\ P\vdash Q \Rightarrow \neg Q \vdash \neg P\\[0.4em]
3.\ P\vdash Q,\ Q\vdash R \Rightarrow P\vdash R\\[0.4em]
4.\ \bot \vdash P\\[0.4em]
5.\ P\vdash N \quad (N\text{ necessary})
\end{array}
\]
\end{whiteboard}

% ============================================================
\section{Rules Later Associated with De Morgan}
% ============================================================

Medieval logicians formulate rules corresponding to what later
became known as De Morgan's laws.

For conjunction:

\[
\neg(P\land Q)
\leftrightarrow
(\neg P\lor \neg Q).
\]

For disjunction:

\[
\neg(P\lor Q)
\leftrightarrow
(\neg P\land \neg Q).
\]

The nineteenth-century name is much later than the logical
insight.

\begin{whiteboard}
\[
\textbf{NEGATION RULES}
\]

\[
\neg(P\land Q)
\leftrightarrow
(\neg P\lor \neg Q)
\]

\[
\neg(P\lor Q)
\leftrightarrow
(\neg P\land \neg Q)
\]

\[
\boxed{
\begin{array}{c}
\text{later name: De Morgan}\\
\text{medieval rules: centuries earlier}
\end{array}
}
\]
\end{whiteboard}

% ============================================================
\section{Primary Logic Re-emerges}
% ============================================================

The medieval theory of consequences contains rules involving:

\begin{itemize}
    \item conjunction,
    \item disjunction,
    \item negation,
    \item conditionals,
    \item contraposition,
    \item transitivity.
\end{itemize}

This is precisely the territory previously identified as primary
or propositional logic.

The tradition therefore reconstructs a substantial portion of the
logical domain earlier developed by the Stoics.

\begin{whiteboard}
\[
\textbf{MEDIEVAL \latin{CONSEQUENTIAE}}
\]

\[
\begin{array}{c}
\neg\\
\land\\
\lor\\
\rightarrow\\
\text{contraposition}\\
\text{transitivity}
\end{array}
\]

\[
\boxed{\text{PRIMARY / PROPOSITIONAL LOGIC}}
\]
\end{whiteboard}

% ============================================================
\section{Beyond Syllogistic: Existential Inference}
% ============================================================

Medieval writers also encounter forms of inference that move
beyond traditional syllogistic.

For example:

\[
\text{Socrates runs.}
\]

Therefore:

\[
\text{Something runs.}
\]

In modern notation:

\[
R(s)
\]

therefore:

\[
\exists x\,R(x).
\]

This is a direct connection between singular predication and
existential quantification.

\begin{whiteboard}
\[
\text{Socrates runs}
\]

\[
\downarrow
\]

\[
\text{Something runs}
\]

\[
\begin{array}{c}
R(s)\\
\hline
\exists x\,R(x)
\end{array}
\]

\[
\boxed{\text{general logic begins to exceed ordinary syllogistic}}
\]
\end{whiteboard}

% ============================================================
\section{The Limits of Medieval Notation}
% ============================================================

Medieval logicians investigate many difficult topics:

\begin{itemize}
    \item relations,
    \item quantification,
    \item scope,
    \item modal contexts,
    \item semantic paradox,
    \item formal replacement,
    \item propositional connectives.
\end{itemize}

But they do not possess a compact symbolic language adequate to
express all of these structures with the clarity later provided by
modern quantificational notation.

The limitation is therefore not simply lack of intelligence or
logical sophistication.

They repeatedly confront problems whose solution requires a
representational advance.

\begin{whiteboard}
\[
\textbf{MEDIEVAL LOGICIANS HAD}
\]

\[
\text{logical insight}
\]

\[
+
\]

\[
\text{technical vocabulary}
\]

\[
\textbf{BUT LACKED}
\]

\[
\boxed{\text{adequate symbolism for quantification}}
\]

\[
\therefore
\]

\[
\text{further progress becomes difficult}
\]
\end{whiteboard}

% ============================================================
\section{Comparison with Chrysippus}
% ============================================================

The later medieval writers can derive one valid consequence from
another and use higher-order principles such as contraposition.

In that respect their work recalls Chrysippus, who derived further
valid schemata from his indemonstrables.

But no surviving medieval work appears to construct a single
highly developed deductive calculus of consequences comparable to
a modern formal system.

\begin{whiteboard}
\[
\textbf{CHRYSIPPUS}
\]

\[
\text{basic schemata}
\]

\[
\downarrow
\]

\[
\text{derived schemata}
\]

\[
\textbf{MEDIEVAL \latin{CONSEQUENTIAE}}
\]

\[
\text{valid consequences}
\]

\[
\downarrow
\]

\[
\text{derive further consequences}
\]

\[
\boxed{\text{but no single unified modern-style calculus}}
\]
\end{whiteboard}

% ============================================================
% SYNTHESIS
% ============================================================

\part{Synthesis}

% ============================================================
\section{What Changed from Rome to the High Middle Ages?}
% ============================================================

The Roman period preserves and combines ancient traditions.

The early Middle Ages recover a restricted logical inheritance.

The twelfth century brings renewed confidence and the work of
Abelard.

The universities then institutionalize logical training and
generate specialized technical traditions.

Two developments become especially characteristic:

\begin{enumerate}
    \item the theory of the properties of terms,
    \item the theory of consequences.
\end{enumerate}

The first pushes toward semantics and quantification.

The second pushes toward a general theory of propositional
inference.

\begin{whiteboard}
\[
\textbf{ROMAN PERIOD}
\]

\[
\text{PRESERVE}
\qquad
\text{TRANSLATE}
\qquad
\text{MIX}
\]

\[
\downarrow
\]

\[
\textbf{EARLY MIDDLE AGES}
\]

\[
\text{RECOVER}
\]

\[
\downarrow
\]

\[
\textbf{ABELARD}
\]

\[
\text{QUESTION}
\qquad
\text{DISTINGUISH}
\qquad
\text{INNOVATE}
\]

\[
\downarrow
\]

\[
\textbf{HIGH MIDDLE AGES}
\]

\[
\begin{array}{c}
\text{SEMANTICS}\\
\text{QUANTIFICATION}\\
\text{PARADOX}\\
\text{CONSEQUENCE}
\end{array}
\]
\end{whiteboard}

% ============================================================
\section{Worked Board Demonstrations}
% ============================================================

\subsection*{1. Relational Inference}

\[
A=2B,
\qquad
C=2A.
\]

Therefore:

\[
C=4B.
\]

The validity depends on relations internal to the propositions.

\begin{whiteboard}
\[
A=2B
\]

\[
C=2A
\]

\[
\therefore C=4B
\]

\[
\boxed{\text{RELATIONAL STRUCTURE matters}}
\]
\end{whiteboard}

% ------------------------------------------------------------

\subsection*{2. The Copula}

Consider:

\[
\text{A chimaera is imaginary.}
\]

The word ``is'' joins subject and predicate.

It does not by itself assert that a chimaera exists as an ordinary
object.

\begin{whiteboard}
\[
\text{A chimaera IS imaginary}
\]

\[
\hspace{0.4in}\uparrow
\]

\[
\text{COPULA}
\]

\[
\boxed{\text{``is'' does not automatically mean ``exists''}}
\]
\end{whiteboard}

% ------------------------------------------------------------

\subsection*{3. Negation and Quantification}

Compare:

\[
\neg(\forall x(M(x)\rightarrow W(x)))
\]

with:

\[
\exists x(M(x)\land \neg W(x)).
\]

In modern logic they are equivalent under the standard
interpretation of the quantifiers.

But Abelard's historical problem concerns the existential
assumptions built into older logical doctrine.

\begin{whiteboard}
\[
\neg[\text{Every }S\text{ is }P]
\]

\[
\text{vs.}
\]

\[
\text{Some }S\text{ is not }P
\]

\[
\boxed{\text{scope + existential assumptions must be examined}}
\]
\end{whiteboard}

% ------------------------------------------------------------

\subsection*{4. Modal Scope}

Compare:

\[
\Diamond\neg\exists x(M(x)\land W(x))
\]

with:

\[
\forall x(M(x)\rightarrow \Diamond\neg W(x)).
\]

The position of the modal operator changes what is asserted.

\begin{whiteboard}
\[
\text{POSSIBLE }[\text{No }S\text{ is }P]
\]

\[
\neq
\]

\[
\text{No }S\text{ is }[\text{POSSIBLY }P]
\]

\[
\boxed{\text{modal scope changes meaning}}
\]
\end{whiteboard}

% ------------------------------------------------------------

\subsection*{5. Paradox of Implication}

Suppose \(P\) is impossible.

Then no case exists in which \(P\) is true and \(Q\) false.

Does that establish:

\[
P\rightarrow Q
\]

for any \(Q\)?

That question separates a truth-conditional conception of
implication from Abelard's stronger demand for intrinsic
connection.

\begin{whiteboard}
\[
P = \text{impossible}
\]

\[
\text{Does }P\rightarrow Q
\]

\[
\text{hold for arbitrary }Q?
\]

\[
\boxed{\text{truth condition vs. intrinsic connection}}
\]
\end{whiteboard}

% ------------------------------------------------------------

\subsection*{6. One Word, Three Suppositions}

Use the same expression in three different contexts.

\[
\text{``man'' is a noun}
\]

\[
\text{man is a species}
\]

\[
\text{every man is mortal}
\]

The logical role of the expression changes.

\begin{whiteboard}
\[
\text{``man''}
\]

\[
\begin{array}{ll}
1. & \text{linguistic sign}\\
2. & \text{species / concept}\\
3. & \text{individual men}
\end{array}
\]

\[
\boxed{\text{same expression --- different supposition}}
\]
\end{whiteboard}

% ------------------------------------------------------------

\subsection*{7. Formal Consequence by Substitution}

Begin with:

\[
\text{Every }M\text{ is }L.
\]

\[
\text{Every }S\text{ is }M.
\]

\[
\therefore \text{Every }S\text{ is }L.
\]

Replace the terms repeatedly.

The validity remains.

\begin{whiteboard}
\[
\begin{array}{c}
\text{Every }M\text{ is }L\\
\text{Every }S\text{ is }M\\
\hline
\text{Every }S\text{ is }L
\end{array}
\]

\[
\boxed{\text{substitution preserves validity}}
\]
\end{whiteboard}

% ------------------------------------------------------------

\subsection*{8. A Small Medieval Rule Set}

A compact group of rules displays how much propositional
structure the medieval tradition had recovered.

\begin{whiteboard}
\[
P\land Q \vdash P
\]

\[
P \vdash P\lor Q
\]

\[
P\rightarrow Q,\ P \vdash Q
\]

\[
P\rightarrow Q,\ \neg Q \vdash \neg P
\]

\[
\neg(P\land Q)
\leftrightarrow
(\neg P\lor\neg Q)
\]

\[
\boxed{\text{substantial PRIMARY LOGIC has re-emerged}}
\]
\end{whiteboard}

% ============================================================
% FINAL SUMMARY
% ============================================================

\section{Final Summary}

The Roman and medieval history of logic should not be described
merely as a period of decline between Greek antiquity and the
Renaissance.

The Roman tradition preserves ancient logical material and makes
it available in Latin.

Boethius becomes the decisive bridge.

The early medieval world inherits only part of ancient logic, but
the twelfth century sees renewed confidence and original logical
work.

Abelard makes major advances in the analysis of propositions,
propositional content, copulation, negation, modality,
quantification, and consequence.

The medieval universities transform logic into a large-scale
academic discipline.

Their interest in sophisms, \latin{insolubilia},
\latin{syncategoremata}, and \latin{obligationes} makes difficult
linguistic and inferential cases into subjects of sustained
technical analysis.

The theory of the \latin{proprietates terminorum} attempts to
explain meaning, reference, and quantificational behavior.

The theory of \latin{consequentiae} broadens the theory of valid
inference and recovers a large part of what can be recognized as
primary or propositional logic.

Medieval thinkers do not yet create modern quantification theory.

A major obstacle is the absence of a sufficiently powerful
symbolic notation.

Nevertheless, they identify many of the problems that modern logic
will later solve in a more systematic symbolic language.

\begin{whiteboard}
\[
\boxed{
\begin{array}{c}
\textbf{ROMAN LOGIC}\\
\text{preserves and transmits}
\end{array}
}
\]

\[
\downarrow
\]

\[
\boxed{
\begin{array}{c}
\textbf{ABELARD}\\
\text{turns recovery into innovation}
\end{array}
}
\]

\[
\downarrow
\]

\[
\boxed{
\begin{array}{c}
\textbf{MEDIEVAL UNIVERSITIES}\\
\text{institutionalize logical analysis}
\end{array}
}
\]

\[
\downarrow
\]

\[
\boxed{
\begin{array}{c}
\textbf{\latin{PROPRIETATES TERMINORUM}}\\
\text{semantics + quantification}
\end{array}
}
\qquad
\boxed{
\begin{array}{c}
\textbf{\latin{CONSEQUENTIAE}}\\
\text{general inference}
\end{array}
}
\]

\[
\downarrow
\]

\[
\boxed{
\begin{array}{c}
\textbf{MEDIEVAL ACHIEVEMENT}\\[0.3em]
\text{A substantial body of the PRIMARY LOGIC}\\
\text{presupposed by syllogistic is articulated.}
\end{array}
}
\]
\end{whiteboard}

% ============================================================
% MASTER WHITEBOARD PACKET
% ============================================================

\section*{Master Whiteboard Packet}

\begin{whiteboard}

\textbf{1. Ancient inheritance}

\[
\text{Aristotle}
+
\text{Chrysippus}
\rightarrow
\text{late ancient mixture}
\rightarrow
\text{Boethius}
\]

\textbf{2. Early medieval recovery}

\[
\text{\latin{ars vetus}}
\rightarrow
\text{trivium}
\rightarrow
\text{universals controversy}
\]

\textbf{3. Abelard}

\[
\text{proposition}
\rightarrow
\text{\latin{dictum}}
\rightarrow
\text{copula}
\]

\[
\text{negation}
\rightarrow
\text{quantification}
\rightarrow
\text{modal scope}
\]

\[
\text{consequence}
\rightarrow
\text{modus ponens / tollens / transitivity}
\]

\textbf{4. University logic}

\[
\text{\latin{ars nova}}
\rightarrow
\text{sophismata}
\rightarrow
\text{\latin{insolubilia}}
\]

\[
\text{\latin{syncategoremata}}
\rightarrow
\text{\latin{obligationes}}
\]

\textbf{5. Term theory}

\[
\text{\latin{significatio}}
\neq
\text{\latin{suppositio}}
\]

\[
\text{material}
\qquad
\text{simple}
\qquad
\text{personal}
\]

\[
\text{distribution}
\rightarrow
\text{quantificational analysis}
\]

\textbf{6. Theory of consequences}

\[
\text{formal}
\qquad
\text{vs.}
\qquad
\text{material}
\]

\[
\text{substitution}
\rightarrow
\text{formal validity}
\]

\[
\neg,\ \land,\ \lor,\ \rightarrow
\]

\[
\boxed{\text{PRIMARY / PROPOSITIONAL LOGIC}}
\]

\textbf{7. Historical limit}

\[
\text{strong logical analysis}
+
\text{weak notation}
\]

\[
\downarrow
\]

\[
\text{difficulty developing full quantification theory}
\]

\textbf{8. Final historical trajectory}

\[
\boxed{
\text{TRANSMISSION}
\rightarrow
\text{RECOVERY}
\rightarrow
\text{INNOVATION}
}
\]

\end{whiteboard}

\end{document}